\documentclass[10pt,a4paper]{amsart}
\usepackage[margin=19mm]{geometry}
\usepackage{amsmath,amssymb,mathtools}
\usepackage[scr=rsfs,cal=euler]{mathalfa}
\usepackage{xfrac}
\usepackage[unicode,hidelinks,bookmarks=false]{hyperref}
\newcommand{\ie}{i.e.\ }
\newcommand{\cf}{cf.\ }
\newcommand{\resp}{respectively\ }
\newcommand{\Subsection}{\subsection}
\providecommand{\nobreakdash}{\nobreak\hbox{-}}
\setlength{\emergencystretch}{2em}
\multlinegap0pt
\usepackage{bm}
\usepackage{bbm}
\usepackage{dsfont}
\usepackage{xspace}
\usepackage{cancel}
\usepackage{mathtools}
\usepackage{amsthm}
\makeatletter
\@ifundefined{theorem}{\newtheorem{theorem}{Theorem}}{}
\@ifundefined{definition}{\newtheorem{definition}[theorem]{Definition}}{}
\@ifundefined{lemma}{\newtheorem{lemma}[theorem]{Lemma}}{}
\@ifundefined{remark}{\newtheorem{remark}[theorem]{Remark}}{}
\makeatother
\newcommand{\Z}{\mathbb{Z}}
\newcommand{\inr}{\operatorname{int}}
\title[Ribbon concordance and fibered predecessors, II: the general]{Ribbon concordance and fibered predecessors, II: the general case}
\author{John A. Baldwin, Jonathan Hanselman, Steven Sivek}
\date{Source: arXiv:2602.21109v1 (CC BY 4.0)}
\begin{document}
\maketitle

\noindent\emph{Source and licence.} Ribbon concordance and fibered predecessors, II: the general case. Version arXiv:2602.21109v1 (CC BY 4.0), distributed under the Creative Commons Attribution 4.0 International licence, as recorded in the arXiv licence metadata for this exact version and shown on the abstract page https://arxiv.org/abs/2602.21109v1. Full rights notice: licenses/arxiv-2602.21109-CC-BY-4.0.txt.\par\smallskip

\noindent\textit{Editorial scope.} Counted unit: the Lemma whose source label is \texttt{lem:winding-number-satellite} in the article (source0.tex lines 527--537, source label \texttt{lem:winding-number-satellite}), delivered together with the article's own complete original proof of it (source0.tex lines 539--564, one \texttt{proof} environment of 26 lines). No mathematics is written, repaired or re-derived: statement and proof are the article's own text line for line. Required context retained uncounted: eq:rational-longitude-vs-muP (lines 467--469); def:generalized-satellite (lines 475--485); def:generalized-companion (lines 493--497); rmk:dual (lines 499--501). Every definition, display and label of those lines is retained unchanged. Omitted from this package and not redistributed: 1--466, 470--474, 486--492, 498--498, 502--526, 538--538, 565--1334. Nothing inside a retained excerpt is omitted or altered: every retained body is delivered exactly as the article's lines are written, so no omission receipt of any kind accompanies this package. Citations: the retained excerpts carry no citation command at all, so no bibliography entry of the article is transcribed and no reference is redistributed. Reference adaptation: none was needed and none was made; every reference inside the delivered unit resolves inside it, so no unresolved reference or citation is produced. Printed slips kept literal: no printed word, symbol, hypothesis or number of the counted statement or of its proof was altered. Layout and class: the article's own class is \texttt{amsart} with packages including amsmath, amssymb, amscd, mathrsfs, url, pinlabel, verbatim, hyperref, geometry, color, dcpic, latexsym, graphicx, epstopdf; the wrapper is the lane's pinned \texttt{amsart} editorial wrapper at 10pt on A4, which restates the theorem-like environments the retained text needs and adds the article's own macro definitions that the retained text needs (\texttt{\textbackslash Z}, \texttt{\textbackslash inr}), with extra wrapper packages (bm, bbm, dsfont, xspace, cancel, mathtools). The delivered numbering is the unit's own. Assets: no retained span contains a figure environment, a graphics-inclusion command, a PDF-page inclusion, a table or a TikZ picture; no image file is copied into this package, the package declares no asset dependency and no image file is redistributed with it. Licence: version arXiv:2602.21109v1 of this article is distributed under the Creative Commons Attribution 4.0 International licence, as recorded in the arXiv licence metadata for this exact version and shown on the abstract page https://arxiv.org/abs/2602.21109v1. A full rights notice is delivered at licenses/arxiv-2602.21109-CC-BY-4.0.txt. Characters: every retained excerpt is pure ASCII, so the delivered engine, XeLaTeX with its default Latin Modern text font, reports no missing character.
\begin{equation} \label{eq:rational-longitude-vs-muP}
[\partial S]=n \alpha_M = nw \cdot [\mu_P]
\end{equation}

\begin{definition}\label{def:generalized-satellite}
A \emph{generalized satellite knot} is a knot $K$ in a rational homology 3-sphere $Y$ whose exterior
\[ E_K = Y \setminus \nu(K) \]
is irreducible, has incompressible boundary, and contains an embedded incompressible torus $T$ that is not boundary-parallel.  We may therefore write
\[ Y = M_1 \cup_T M_2, \qquad \partial M_1 = T = -\partial M_2,\qquad K \subset \inr(M_2). \]
The fact that $H_2(Y;\Z)=0$ implies via the Mayer--Vietoris sequence that
\begin{equation} \label{eq:h2-satellite}
H_2(M_1;\Z) = H_2(M_2;\Z) = 0,
\end{equation}
and we further require as part of this definition that the rational longitudes $\alpha_{M_1}, \alpha_{M_2} \subset T$ have pairwise distance one. We  refer to $T$ as the \emph{satellite torus.}
\end{definition}

\begin{definition} \label{def:generalized-companion}
Let $\mu_C \subset \partial M_1$ be the curve identified with $\alpha_{M_2} \subset \partial M_2$ in $T$, and let
\[ Z = M_1(\mu_C) \]
be the closed 3-manifold obtained by Dehn filling $M_1$ along $\mu_C$.  The \emph{companion knot} \[C \subset Z\] is given by the core of the filling solid torus.
\end{definition}

\begin{remark}\label{rmk:dual}
The companion $C$ has exterior $M_1$, and in the torus $T = \partial M_1$ its rational longitude $\lambda_C = \alpha_{M_1}$ is geometrically dual to its meridian $\mu_C = \alpha_{M_2}$, by Definition \ref{def:generalized-satellite}.
\end{remark}

\begin{lemma} \label{lem:winding-number-satellite}
Let $K \subset Y$ be a nullhomologous generalized satellite knot, with satellite torus $T \subset Y$ and corresponding decomposition \[Y = M_1 \cup_T M_2.\]  Let $P \subset M_2$ and $C \subset Z$ be the associated pattern and companion knots.  Let \[E_K = Y \setminus \nu(K), \qquad E_P = M_2\cap E_K = M_2 \setminus \nu(P),\] with $\mu_K$ and $\lambda_K$ denoting the meridian and Seifert longitude of $K$ in each exterior.
If $P$ has winding number $w \geq 0$, then $w$ is an integer and any Seifert surface $\Sigma$ for $K$ meeting $T$ transversely satisfies
\[ [\Sigma \cap T] = w[\lambda_C] \]
as elements of $H_1(T;\Z)$.  In particular, this gives us relations $[P] = w[\lambda_C]$ in $H_1(M_2;\Z)$, and
\begin{align*}
[\lambda_K] &= w[\lambda_C], \\
nw[\mu_K] &= n[\mu_C]
\end{align*}
in $H_1(E_P;\Z)$, where $n \geq 1$ is the order of the rational longitude $\alpha_{M_2}$.
\end{lemma}

\begin{proof}
Let $\Sigma \subset E_K$ be a Seifert surface for $K$ which is transverse to $T$.  Let $k \geq 1$ denote the order of the rational longitude $[\lambda_C] = \alpha_{M_1}$ in $H_1(M_1;\Z)$. 
Observe that $\Sigma \cap M_1$ is a properly embedded surface in $M_1$, so its boundary class
\begin{equation*} \label{eq:boundary-sigma-y0}
\partial_*[\Sigma \cap M_1] \in H_1(\partial M_1;\Z) = H_1(T;\Z)
\end{equation*}
is in the kernel of the inclusion-induced map \[H_1(T;\Z)\to H_1(M_1;\Z),\]  and is thus an integer multiple
\begin{equation}\label{eq:boundary-mk}
\partial_*[\Sigma \cap M_1]= km[\lambda_C]\end{equation}
of $k$ times the rational longitude.  Note that \[\Sigma \cap  E_P\subset M_2\] extends to the core of the tubular neighborhood $\nu(P)$ to have boundary $P$, producing a relation
\begin{equation} \label{eq:P-lambdaC}
[P] = km[\lambda_C]
\end{equation}
in $H_1(M_2;\Z)$.

Now, the meridian $\mu_C$ and rational longitude $\lambda_C$ are  dual to each other in $T$, as noted in Remark \ref{rmk:dual}, and $\mu_C$ is identified in $T$ with the rational longitude $\alpha_{M_2}$ of $M_2$, by Definition \ref{def:generalized-companion}.  Letting $S$ be the surface generating 
\[H_2(M_2,\partial M_2;\Z) \cong \Z,\] with $\partial_*[S] = n\cdot \alpha_{M_2}$ and $n \geq 1$, we deduce that the intersection pairing on $M_2$ then satisfies
\begin{equation} \label{eq:S-lambdaC}
[S] \cdot [\lambda_C] = n,
\end{equation}
as can be seen by pushing $\lambda_C$ slightly into $M_2$ so that it meets $S$ in a single point near each copy of $\alpha_{M_2}$ in $\partial S$.  Applying \eqref{eq:S-lambdaC} and then \eqref{eq:P-lambdaC}, we have
\[ kmn = km\left([S] \cdot [\lambda_C]\right) = [S] \cdot km[\lambda_C] = [S] \cdot [P] = nw, \]
so finally $w = km$.  In particular $w \in \Z$, so the relations \eqref{eq:boundary-mk} and \eqref{eq:P-lambdaC} become
\[ [\Sigma \cap T] = \partial_*[\Sigma \cap M_1] = km [\lambda_C]=w[\lambda_C] \]
in $H_1(T;\Z)$ and \[[P] = w[\lambda_C]\] in $H_1(M_2;\Z)$.  Similarly, the intersection $\Sigma \cap E_P \subset M_2$  gives rise to the relation \[[\lambda_K] = w[\lambda_C]\] in $H_1(E_P;\Z)$.  Finally, the relation \[nw[\mu_K] = [\mu_C]\] is precisely \eqref{eq:rational-longitude-vs-muP} after we identify $\mu_K = \mu_P$ and $\mu_C = \alpha_{M_2}$.
\end{proof}

\end{document}
